\documentclass[11pt]{article}
\usepackage[a4paper,margin=1in]{geometry}
\usepackage{amsmath,amssymb,mathtools,bm}
\usepackage{enumitem,booktabs,array}
\usepackage{tcolorbox}
\usepackage{hyperref}
\usepackage[protrusion=true,expansion=false]{microtype}
\hypersetup{colorlinks=true,linkcolor=blue,urlcolor=blue,citecolor=blue}
\setlist[itemize]{topsep=2pt,itemsep=2pt}
\setlist[enumerate]{topsep=2pt,itemsep=2pt}
\newtcolorbox{keybox}{colback=blue!3!white,colframe=blue!40!black,boxrule=0.4pt,arc=1mm}
\newtcolorbox{alertbox}{colback=red!3!white,colframe=red!40!black,boxrule=0.4pt,arc=1mm}
\newtcolorbox{answerbox}{colback=green!3!white,colframe=green!40!black,boxrule=0.4pt,arc=1mm}
\newtcolorbox{drillbox}{colback=yellow!5!white,colframe=orange!60!black,boxrule=0.4pt,arc=1mm}
\newcommand{\EV}{\mathbb{E}}
\newcommand{\Var}{\mathrm{Var}}
\newcommand{\Cov}{\mathrm{Cov}}
\newcommand{\blank}[1]{\underline{\hspace{#1}}}

\title{W04-04: Ivashina \& Scharfstein (2010, JFE) \\
\large ``Bank Lending During the Financial Crisis of 2008'' --- Active Recall Workbook}
\author{Banking \& Risk Management --- Exam Prep}
\date{\today}
\begin{document}\maketitle

\begin{keybox}
\textbf{Paper in one sentence.} Using Dealscan syndicated-lending data, IS document that new loan origination collapsed by $\sim 47\%$ in the peak of the 2008 crisis (Q4 2008 vs.\ Q2 2007), and that \emph{banks with more exposure to Lehman co-syndications} and \emph{less stable deposit funding} cut lending the most --- pinning the collapse on bank balance-sheet stress rather than solely on borrower demand.

\textbf{Placement.} The textbook ``lending during the 2008 crisis'' paper. Complements Gorton--Metrick (W04-01) on the wholesale-funding shock and Mian--Sufi (W04-03) on household credit. Central reference for bank-supply effects in the Great Recession.
\end{keybox}

\section*{Round 1: Complete Walk-Through}

\subsection*{1.1 Context}
The 2008 crisis culminated with Lehman's bankruptcy (15 Sept 2008) and a broad collapse in wholesale funding: ABCP, repo, and interbank markets froze. Bank balance sheets were strained by falling asset values and rising liquidity demand from drawn credit lines.

\subsection*{1.2 Data}
Dealscan loan-level syndicated-lending data: deal origination, participants, lead bank, size, purpose. Quarterly aggregates Q3 2007--Q4 2008.

\subsection*{1.3 Aggregate facts}
\begin{itemize}
\item New loan volume falls steadily through 2008 and \emph{collapses} 47\% in Q4 2008 vs.\ the Q2 2007 peak.
\item Investment-grade and non-investment-grade lending both fall, though the composition tilts more conservatively.
\item Existing \emph{drawn} credit lines spike (firms drawing down as liquidity insurance).
\end{itemize}

\subsection*{1.4 Cross-sectional identification}
Two main tests:
\begin{enumerate}
\item \textbf{Lehman-co-syndication exposure.} Banks that had co-syndicated more with Lehman were arguably more affected by Lehman's collapse. These banks cut new origination significantly more after September 2008.
\item \textbf{Deposit vs.\ wholesale funding mix.} Banks with more stable deposit funding cut lending less; banks reliant on wholesale funding cut more. (An early test of the deposits-channel intuition later formalised by DSS.)
\end{enumerate}

\subsection*{1.5 Drawn-line channel}
Firms with existing credit lines drew them down during the crisis to hoard liquidity. This raised \emph{actual} bank lending even as origination collapsed --- a distinct margin that IS emphasise.

\subsection*{1.6 Caveats}
\begin{alertbox}
\begin{itemize}
\item Syndicated lending is only a slice of total bank credit; SME lending is not directly observed.
\item Selection: only firms that successfully close a deal appear; unmet demand is unobserved.
\item Demand shifts during the crisis (firms cutting investment) could account for part of the volume decline. IS handle this by focusing on cross-bank variation.
\end{itemize}
\end{alertbox}

\section*{Round 2: Level 1 Fill-in}
\begin{enumerate}
\item Data source: \blank{2cm} syndicated-loan data.
\item Volume collapse: $\sim$\blank{1cm}\% Q4 2008 vs.\ Q2 2007.
\item Cross-bank treatment 1: \blank{2cm} co-syndication exposure.
\item Cross-bank treatment 2: reliance on \blank{2cm} funding (vs.\ deposits).
\item Drawn-line effect: \emph{actual} lending can \blank{1cm} as origination falls.
\end{enumerate}

\section*{Round 3: Level 2 Prompts}
\begin{enumerate}
\item Why does comparing Lehman-exposed to non-Lehman-exposed banks help identify a supply channel?
\item Why do firms rush to draw credit lines during a crisis? What does this reveal about bank balance-sheet dynamics?
\item How does IS relate to KRS (2002) on the synergy of deposits and credit lines?
\item What data would you need to test for the same effects at the SME-lending margin?
\end{enumerate}

\section*{Round 4: Minimal Scaffolding}
From a blank page: (i) data, (ii) aggregate collapse, (iii) two cross-bank tests, (iv) drawn-line channel, (v) caveats.

\section*{Round 5: Blank Page}
Essay: the 2008 lending collapse as a bank-supply shock.

\vspace{6pt}
\begin{drillbox}
\textbf{Speed Drill.} (1) Dataset. (2) Peak-to-trough origination decline. (3) Lehman exposure test. (4) Funding-mix test. (5) What drawn lines do to realised credit.
\end{drillbox}

\section*{Quick Reference \& Answer Key}
\begin{answerbox}
\begin{itemize}
\item \textbf{Aggregate.} New syndicated lending $\downarrow$ $\sim 47\%$ Q2 07 $\to$ Q4 08.
\item \textbf{Supply evidence.} Lehman-exposed banks and wholesale-funded banks cut more.
\item \textbf{Drawn lines.} Firms hoard liquidity by drawing existing commitments.
\item \textbf{Lesson.} Bank-balance-sheet stress is a first-order driver of the lending collapse.
\end{itemize}
\end{answerbox}

\begin{keybox}
\textbf{30-second exam pitch.} Ivashina \& Scharfstein (2010) document the 2008 lending collapse using Dealscan syndicated-lending data. New-loan origination falls by nearly half from the Q2 2007 peak to Q4 2008. The decline is sharpest at banks with two signs of balance-sheet stress: (i) more co-syndication exposure to Lehman, and (ii) more reliance on wholesale (non-deposit) funding. Meanwhile, firms drew down existing credit lines en masse, so realised bank credit could rise even as origination cratered --- a channel that disciplines naive reading of aggregate bank credit during a crisis. The paper is the bank-supply bookend of the 2008 story: Gorton--Metrick on the repo run, Mian--Sufi on household credit, and IS on the corresponding contraction in corporate credit from stressed banks.
\end{keybox}

\end{document}
